\section{Resummation of the heat kernel expansion}\label{sec8}
The heat kernel coefficients define the one-loop counterterms
in the background field formalism. In many cases the heat kernel
can also give a useful information on the finite part of the
effective action. Just one of the examples is the large mass
expansion (\ref{largemass}) which is valid when all background
fields and their derivatives are small compared to the mass of
the quantum field. In order to get the effective action in other
limiting cases one has to re-arrange the heat kernel 
expansion\footnote{A diagrammatic technique which can be used in resummations
of the heat kernel expansion is described in \cite{Moss:1999wq}.}.
%%%%%%
\subsection{Modified large mass expansion}\label{s8finite}
In many physical applications there is a quantity $\mathcal{M}^2$
which is large compared to the rest of the background fields
and their derivatives. Therefore, it is a well motivated problem
to construct an expansion of the effective action in a
power series of $\mathcal{M}^{-1}$. To do so, one has to re-express
the heat kernel as
\begin{equation}
K(t;D)=e^{-t\mathcal{M}^2} \sum_k t^{(k-n)/2} \tilde a_k \,.
\label{modifHK}
\end{equation}
To obtain the effective action one has to integrate the heat kernel
over $t$ (cf eq.\ (\ref{WKtD})). To simplify the argumentation
we assume the cut-off regularization (\ref{WLamreg}), although in other
regularization schemes the result will be essentially the same.
The divergent and finite parts of the effective action are given
by (\ref{divWLam}) and (\ref{largemass}) respectively up to the obvious
replacements: $m\to \mathcal{M}$, $(4\pi )^{-n/2}b_k \to \tilde a_k$.
After a suitable renormalization, this procedure indeed gives a large
$\mathcal{M}$ expansion of the effective action.

Therefore, the problem is reduced to calculation of 
$\tilde a_k$. If $\mathcal{M}$ commutes with $D$ the coefficient
$\tilde a_k$ are the heat kernel coefficients for
the operator $D-\mathcal{M}^2$. This case returns us to the
standard large mass expansion. If $\mathcal{M}^2$ does not commute
with $D$ then
\begin{equation}
e^{-tD}\ne e^{-t(D-\mathcal{M}^2)}e^{-t\mathcal{M}^2} \,.\label{noteqexps}
\end{equation}
To achieve an equality the right hand side of (\ref{noteqexps}) must
be corrected by commutator terms. In this case calculation of
$\tilde a_k$ requires some amounts of extra work.

Heavy particles of non-equal masses described by the mass matrix
$\mathcal{M}^2$ is probably the most immediate example of a physical
system to which the modified large mass expansion should be applied.
Corresponding technical tools were developed recently 
\cite{Osipov:2000rg,Osipov:2001bj,Osipov:2001nx}.

The next example is a scalar field in curved space (cf. sec\ 
\ref{s3scalar}). Parker and Toms suggested 
\cite{Parker:1985dj,Parker:1985mv} to use the modified large
mass expansion with $\mathcal{M}^2=m^2+\xi R$ which partially
sums up contributions of the scalar curvature $R$ to the effective
action. This formalism was developed further in 
\cite{Jack:1985mw,Sushkov:2002rf}.

In some cases the term-by-term integration of the heat kernel
expansion gives good estimates of the vacuum energy even if no
large parameter is explicitly present in the model (see, e.g.
calculations of quantum corrections to the mass of two dimensional
solitons \cite{Dunne:1999du,AlonsoIzquierdo:2002eb}). 
%%%%%%
\subsection{Covariant perturbation theory}\label{s8covpt}
Suppose that the matrix potential $E$, the Riemann curvature,
and the field strength $\Omega$ are small but rapidly varying.
Can we get any information on the heat kernel coefficients
containing a fixed power of the quantities listed above and
{\it arbitrary} number of derivatives? For the linear and
quadratic orders on a manifold without boundaries the answer
is positive and the results can be obtained either by the
functorial methods of sec.\ \ref{s4gen} 
\cite{Gilkey:1988,Branson:1990b}
or by solving the DeWitt equation 
\cite{Avramidi:1989er,Avramidi:1990ug,Avramidi:1991je}
(cf. sec.\ \ref{s4iter}). Retaining only the leading terms
we have for $k\ge 1$:
\begin{eqnarray}
&&a_{2k}(x,D) =(4\pi )^{-n/2} \left[ \alpha_1(k) 
\Delta^{k-1} R +\alpha_2 (k) \Delta^{k-1} E \right.\label{leadhk}\\
&&\qquad\qquad\qquad\qquad \left. 
+\mbox{higher order terms} \right]\,, \nonumber
\end{eqnarray}
where \cite{Gilkey:1979,Gilkey:1988}
\begin{equation}
\alpha_1(k)=\frac{k!k}{(2k+1)!} \,,\qquad
\alpha_2(k)=\frac{2k!}{(2k)!}\,.\label{al12k}
\end{equation}

One can sum up the expansion (\ref{leadhk}) and corresponding higher
order terms to obtain an information on the behaviour of the full
heat kernel in the limit described above. There is, however, a more
straightforward method to control non-localities which is called
the covariant perturbation theory 
\cite{Barvinsky:1987uw,Barvinsky:1990up,Barvinsky:1990uq}. To make the
idea of this method most transparent we consider a simple example
of flat $M$ and zero connection $\omega$. The exponent
$\exp (-tD)=\exp (t(\Delta + E))$ with $\Delta =\partial_\mu^2$
can be expanded in a power series
in $E$:
\begin{equation}
e^{-tD}=e^{t\Delta} + \int\limits_0^t ds e^{(t-s)\Delta} E e^{s\Delta}
+\int\limits_0^t ds_2 \int\limits_0^{s_2} ds_1
e^{(t-s_2)\Delta}E e^{(s_2-s_1)\Delta} E e^{s_1\Delta} +\dots
\label{pertex}
\end{equation}
The heat trace can be also expanded,
\begin{equation}
K(t,D)={\rm Tr}\left( e^{-tD}\right)=
\sum_{j=0}^\infty K_j(t),\label{pexhtr}
\end{equation}
where $K_j$ contains the $j$th power of $E$.

Covariant perturbation theory approach prescribes to take the
$0$th order heat kernel in the free space form (cf. eq.\ 
(\ref{simplestHK}) for $m=0$):
\begin{equation}
K_0(x,y;t)=(4\pi t)^{-n/2} \exp
\left( -\frac{(x-y)^2}{4t} \right) \,,\label{Kzero}
\end{equation}
which is an exact kernel on $M=\mathbb{R}^n$ only\footnote{The
original paper \cite{Barvinsky:1987uw} contained a curved space
generalisation of (\ref{Kzero}). This, however, does not improve
the global issues discussed below.}. This formula neglects all
global contributions and, therefore, is valid only for sufficiently
close $x$ and $y$ and for small $t$. We have:
\begin{equation}
K_0(t)=\iM {\rm tr}_V K_0(x,x;t)=(4\pi t)^{-n/2} \iM {\rm tr}_V (I_V)
\,,\label{trKzero}
\end{equation}
where $\ptr I_V$ simply counts discrete (spin and internal) 
indices of $D$. This formula reproduces the $a_0$ contribution
to the heat kernel.

In the next order, we have:
\begin{eqnarray}
K_1(t)=&&{\rm Tr} \left( \int_0^t ds e^{(t-s)\Delta}Ee^{s\Delta}
\right)=
{\rm Tr}\left( \int_0^t ds e^{t \Delta} E \right) \nonumber \\
=&& t\, \ptr \iM K_0 (x,x;t) E(x) \nonumber\\
=&&\frac t{(4\pi t)^{n/2}}
\iM \ptr E(x) \,,\label{Kone}
\end{eqnarray}
where we used the cyclic property of the trace and the expression
(\ref{Kzero}) for the $0$th order heat kernel. This expression is
consistent with the $k=1$ term of (\ref{leadhk}). The terms with
$k>1$ are total derivatives, and, therefore, they do not contribute
to the integrated heat kernel. 

The quadratic order of the heat kernel reads
\begin{eqnarray}
K_2(t)=&&{\rm Tr} \left( \int\limits_0^t ds_2 \int\limits_0^{s_2} ds_1
e^{(t-s_2)\Delta}E e^{(s_2-s_1)\Delta} E e^{s_1\Delta} \right) \nonumber \\
=&&\ptr \int_M dy \int_M dz 
\int\limits_0^t ds_2 \int\limits_0^{s_2} ds_1 
K_0 (z,y;t-s_2+s_1)E(y)\nonumber\\
&&\qquad\qquad\qquad\qquad\qquad \times K_0(y,z;s_2-s_1)E(z) \,.\label{K2a}
\end{eqnarray}
Next we introduce the rescaled variables $\xi =s/t$ and get rid of redundant
integrations to obtain:
\begin{eqnarray}
K_2(t)=&&\frac {t^2}2 \ptr \int_M dy \int_M dz \int\limits_0^1 d\xi
K_0(z,y;t(1-\xi )) E(z) \nonumber\\
&&\qquad\qquad\qquad\qquad\qquad 
\times K_0(y,z;t\xi)E(z) \,.\label{K2b}
\end{eqnarray}
Now we use the identity
\begin{equation}
K_0(z,y;t(1-\xi )) K_0(y,z;t\xi)=(4\pi t)^{-n/2}
K_0(z,y;t\xi(1-\xi)) \label{KKKide}
\end{equation}
and relate the heat kernel on the right hand side of (\ref{KKKide})
to a matrix element of $\exp (t\xi(1-\xi )\Delta )$ to obtain
the final result:
\begin{equation}
K_2(t)=\frac{t^2}{(4\pi t)^{n/2}}
\iM E f(-t\Delta ) E \,,\label{Ktwo}
\end{equation}
where the non-local form-factor $f$ reads:
\begin{equation}
f(q)=\frac 12 \int\limits_0^1 d\xi e^{-q\xi (1-\xi)}=
\frac 12 e^{-q/4} \sqrt{\pi /q} {\rm Erfi} \left[ \sqrt{q}/2 \right]\,.
\end{equation}
As we have already discussed above, applicability of this formula is
limited by our choice of $K_0(x,y;t)$. Namely, the potential $E(x)$
should have a compact support and $t$ should be reasonably small.
Note, that the small $t$ approximation in the expansion (\ref{pertex})
is self-consistent: if $t$ is small, the integration variables $s_i$
are even smaller.

The main difficulty in constructing an expansion in powers
the Riemann curvature and of the field strength is to organise
the procedure in a {\it covariant} way. The details of the
construction and higher order form-factors can be found in  
\cite{Barvinsky:1987uw,Barvinsky:1990up,Barvinsky:1990uq,Barvinsky:1994hw,
Barvinsky:1994ic}. From further developments of this method
we mention the work of Gusev and Zelnikov \cite{Gusev:1999cv}
who demonstrated that in two dimension one can achieve considerable
simplifications in the perturbation expansion by using the dilaton
parametrisation of the potential (cf. eq.\ (\ref{ehnEh})).
Recently Barvinsky and Mukhanov \cite{Barvinsky:2002uf} suggested
a new method for calculation of the non-local part of the effective
action based on the resummation of the perturbation series for the
heat kernel. This method was extended in \cite{Barvinsky:2003rx} to
include late time asymptotics of the heat kernel in curved space.

Let us stress that the expansion (\ref{pertex}) can be used also
for singular potentials. For example, it is very effective for 
the calculation of the heat kernel for $\delta$-potentials
concentrated on a co-dimension one subsurface 
\cite{Gaveau:1986,Bordag:1999ed,Moss:2000gv} (cf.\ sec.\ \ref{s6dom}).
If the $\delta$-potential has its support on a submanifold
of co-dimension greater than one, the expansion diverges 
\cite{Bordag:1999ed}. 
With a suitable choice of the zeroth order heat kernel $K_0(x,y;t)$
and of an operator to replace $E$ in (\ref{pertex}) one can treat
manifolds with boundaries \cite{Bordag:2001fj} where the perturbative
expansion takes the form of the multiple reflection expansion
\cite{Balian:1970,Balian:1977za,Hansson:1983xt}.
%%%%%%
\subsection{``Low energy'' expansion}\label{s8low}
Let us now turn to the opposite case when the derivatives are less
important when the potential and the curvatures. In this case
one has to collect the terms which are of a fixed order in the
derivatives, but contain all powers of $E$, $R$ and $\Omega$.
Since the derivatives are sometimes identified with the energy,
this approximation is being called the low energy expansion.
This scheme goes back to Schwinger's calculations
\cite{Schwinger:1951} in constant electromagnetic 
fields\footnote{We like to mention also the paper
\cite{Heisenberg:1936} which treated the effective action
in external electromagnetic field from a different point view.}.

Let us consider a simple example \cite{Blau:1991iz}. Let
$M=\mathbb{R}^2$. Consider a scalar particle in constant
electromagnetic field with the field strength $F_{12}=-F_{21}=B$.
As a potential we choose $A_1=0$, $A_2=Bx^1$. Then the operator
acting on quantum fluctuations is
\begin{equation}
D=-\partial_1^2-(\partial_2 -iBx^1)^2 +m^2 \,.\label{cmagD}
\end{equation}
This operator commutes with $\partial_2$. Therefore, we can look
for the eigenfunctions of $D$ in the form 
$\phi_k (x)=\exp (ikx^2) \tilde\phi (x^1)$. 
\begin{equation}
D\phi_k (x) = \left( -\partial_1^2 +B^2 (x^1 - k/B)^2 +m^2 \right)
\phi_k (x) \,.\label{cmDk}
\end{equation}
In the $x^1$ direction we have the one-dimensional harmonic
oscillator potential (cf. (\ref{harmos})). Therefore, the 
eigenvalues are $\lambda_{k,p}= (2p+1)|B|+m^2$. These eigenvalues
do not depend on $k$. For this reason the heat kernel
\begin{equation}
K(t,D)=\sum_{k,p} \exp (-t\lambda_{k,p}) \label{cmagKt}
\end{equation}
is ill defined. To overcome this difficulty it was suggested 
\cite{Blau:1991iz} to put the system in a box (without specifying
any boundary conditions, however) and replace the sum over $k$
by the degeneracy factor $(|B| {\rm Vol})/(2\pi )$, where ${\rm Vol}$
is volume of the box. The degeneracy factor is chosen in such a way
that the resulting heat kernel
\begin{equation}
K(t,D)= \frac{B{\rm Vol}}{4\pi } e^{-m^2t} [{\rm sinh} (Bt)]^{-1}
\label{cmagKB}
\end{equation}
behaves as ${\rm Vol}/(4\pi t)$ in the limit $B\to 0$.

Calculations of the effective action in covariantly constant
background gauge fields by the spectral theory methods have 
been performed by many authors 
\cite{Batalin:1977uv,Dittrich:1983ej,Elizalde:1984yb,
Elizalde:1985zv,Blau:1991iz}. Such calculations were motivated, at least
partially, by various models of quark confinement (see 
\cite{Pagels:1978dd,Adler:1982rk} and references therein).

There exists an algebraic method \cite{Avramidi:1993yf,Avramidi:1995ik} 
which allows to evaluate the low energy heat kernel by using 
exclusively the commutator algebra. Let us briefly formulate the
results of \cite{Avramidi:1993yf,Avramidi:1995ik}. Consider a flat
manifold ($R_{\mu\nu\rho\sigma}=0$) with the background fields
satisfying the ``low-energy conditions'':
\begin{equation}
\nabla_\mu \Omega_{\mu\nu}=0,\qquad
\nabla_\mu\nabla_\nu\nabla_\rho E=0 
\label{lowecon}
\end{equation}
with usual definitions of $E$ and $\Omega$ (see (\ref{D1}) - 
(\ref{Omega})). Moreover, let us suppose that the background
is ``approximately abelian'', i.e. that $\Omega_{\mu\nu}$,
$E$ and all their covariant derivatives are mutually commuting matrices. 

With all these assumptions a closed expression for the heat kernel
may be obtained:
\begin{equation}
K(t;x,x;D)=(4\pi t)^{-n/2} \exp \left(
tE +\Phi (t) +\frac {t^3}4 E_{;\mu} \Psi^{\mu\nu} E_{;\nu} \right)\,,
\label{avrahk}
\end{equation}
where $\Phi$ and $\Psi$ are complicated functions of $t$, $E$ and $\Omega$
\cite{Avramidi:1995ik}. If $\nabla_\mu\nabla_\nu E=0$,
\begin{eqnarray}
&&\Phi (t)=-\frac 12 \ln \det \left( \frac{\sinh t \Omega}{t\Omega}
\right) \,,\nonumber\\
&&\Psi (t) = (t\Omega )^{-2} \left( t\Omega \coth (t\Omega )-1
\right) \,,\label{PhiPsi1}
\end{eqnarray}
where $\Omega$ has to be understood as a space-time matrix
$\Omega_{\mu\nu}$, so that multiplication in (\ref{PhiPsi1})
is the matrix multiplication, and $\det$ is the determinant
of an $n\times n$ matrix. These formulae generalise the equation
(\ref{cmagKB}) and justify the choice of the degeneracy factor made
to derive it.

If $\Omega =0$,
\begin{eqnarray}
&&\Phi (t) = -\frac 12 \ln \det \left( 
\frac{\sinh (2t\sqrt{P}}{2t\sqrt{P}} \right) \,,\nonumber\\
&&\Psi (t) = -\left( t\sqrt{P} \right)^{-3}
\left( \tanh (t\sqrt{P}) -t\sqrt{P}) \right) \,,\label{PhiPsi2}
\end{eqnarray}
where $P_{\mu\nu}=-(1/2)\nabla_\mu\nabla_\nu E$. In the particular
case of one-dimensional harmonic oscillator these formulae
reproduce (\ref{harmheat}). 

In curved space the best one can do in the framework of the
low-energy expansion is to consider
locally symmetric manifolds, i.e. vanishing $\Omega$,
covariantly constant Riemann curvature and constant $E$. 
In this case formulae similar to the ones presented above
are available \cite{Avramidi:1994zc,Avramidi:1996zp}.
Covariantly constant curvature means that locally 
the manifold $M$ is a symmetric space. Various approaches
to the heat kernel on such manifolds are described in detail
in monographs and survey articles 
\cite{Hurt:1983,Camporesi:1990wm,Bytsenko:2001xs}. In particular,
very detailed information may be obtained for group manifolds
(see, for example, \cite{Dowker:1971vu}) and for hyperbolic
spaces \cite{Bytsenko:1996bc}.
%%%%%%
\subsection{Heat kernel on homogeneous spaces}\label{s8homo}
In this section we briefly explain how one can find
the spectrum of some ``natural'' differential
operators on homogeneous spaces
by purely algebraic methods. We start with
some basic facts from differential geometry and harmonic analysis
\cite{Kobayashi:1969,Helgason:1984}. 
Consider a homogeneous space $G/H$ of
two compact finite-dimensional Lie groups $G$ and $H$.
The Lie algebra $\mathcal{G}$ of $G$ can be decomposed as
\begin{equation}
\mathcal{G}=\mathcal{H}\oplus \mathcal{M} \,,\label{Gdecom}
\end{equation}
where $\mathcal{H}$ is the Lie algebra of $H$ and $\mathcal{M}$
is the complement of $\mathcal{H}$ in $\mathcal{G}$ with
respect to some bi-invariant metric. We have:
\begin{equation}
[\mathcal{H},\mathcal{M}] \subset \mathcal{M} \,,\label{HMM}
\end{equation}
where $[\ ,\ ]$ is the Lie bracket on $\mathcal{G}$. If, moreover,
$[\mathcal{M},\mathcal{M}]\subset \mathcal{H}$, then $G/H$ is a
symmetric space. We do not impose this restriction.

$\mathcal{M}$ can be identified with the tangent space to $M=G/H$
at the origin (i.e. at the point which represents the unit element of $G$).
Eq.\ (\ref{HMM}) tells us that $\mathcal{H}$ acts on $\mathcal{M}$
by some (orthogonal) representation. This action defines the embedding
\begin{equation}
\mathcal{H}\subset so(n) \,.\label{Hinso}
\end{equation}
All physical fields are classified according to certain representations
of the Lie algebra $so(n)$. Restrictions of these representations to
$\mathcal{H}$ define transformation properties of the field with
respect to $\mathcal{H}$. From now on we work with each irreducible
representation of $\mathcal{H}$ separately.

The field $\Phi^A$ belonging to
an irreducible representation $\mathcal{T}(H)$ can be expanded as 
(see, e.g., \cite{Salam:1982xd}):
\begin{equation}
 \Phi_A (x)=({\rm Vol})^{-\frac 12} \sum_{j, \xi ,q}
\sqrt {\frac {d_j}{d_{\mathcal{T}}}} \mathcal{T}^{(j)}_{A\xi ,q}
(g_x^{-1}) \phi^{(j)}_{q,\xi }\ , \label{eq:(A.2)}
\end{equation}
where ${\rm Vol}$ is the volume of $G/H$, 
$d_{\mathcal{T}}={\rm dim}\mathcal{T}(H)$.
We sum over the representations $\mathcal{T}^{(j)}$ of $G$
which give $\mathcal{T}(H)$ after reduction to $H$. $\xi$ labels
multiple components  $\mathcal{T}(H)$ in the branching
$\mathcal{T}^{(j)}\downarrow H$, $d_j=\dim \mathcal{T}^{(j)}$. 
$q$ runs from $1$ through $d_j$. The matrix
elements of $\mathcal{T}^{(j)}$ have the following orthogonality
property
\begin{equation} \iM
\mathcal{T}^{(j) \dag}_{A\zeta ,q} (g_x^{-1})
\mathcal{T}^{(j')}_{A\xi ,p} (g_x^{-1})
=({\rm Vol})d_j^{-1} d_{\mathcal{T}} \delta_{\zeta \xi} \delta_{p q}
\delta_{jj'} \label{eq:(A.3)}\end{equation}
Therefore, to construct the harmonic expansion on $G/H$ it is
necessary to have powerful methods for reduction of the representations 
from $G$ to $H$. There are several standard 
\cite{Barut:1977,Slansky:1981yr}
and less standard  \cite{Lyakhovsky:1989cr,Lyakhovsky:1991vv}
techniques which may be used depending on the particular
homogeneous space.

It is important that not only the harmonic expansion but also
the spectrum of the invariant operators can be analysed by the group
theoretical methods. The covariant derivative on $G/H$
reads \cite{Kobayashi:1969}:
\begin{equation}
\nabla_\mu = \nabla_\mu^{[c]}+\Gamma^{[R]}_\mu \,.\label{hscovdiv}
\end{equation}
Here $\nabla^{[c]}$ is the canonical covariant derivative on
$G/H$. At the origin $\nabla^{[c]}$ can be identified with 
the tangent space generators from $\mathcal{M}$ taken in
the representation $\mathcal{T}^{(j)}$. The part $\Gamma^{[R]}_\mu$
depends on the invariant metric on $G/H$ and on the
structure constants of $\mathcal{G}$ restricted to $\mathcal{M}$.
On symmetric spaces such structure constants are zero and,
therefore, the Laplace-Beltrami operator has a particularly
simple form:
\begin{equation}
D\simeq \nabla^{\mu[c]}\nabla_\mu^{[c]} \simeq C_2 (G)-C_2 (H)\,,
\label{LapCas}
\end{equation}
where $C_2$ are quadratic Casimir operators of $G$ and $H$ which
depend on the representations $\mathcal{T}^{(i)}$ and 
$\mathcal{T}(H)$ respectively.
On general homogeneous spaces the expressions are a bit more
complicated (see, e.g., 
\cite{Vassilevich:1988em,Lyakhovsky:1991vv} for explicit
examples). In any case, eigenvalues of $D$ are given by
a second order polynomial $Q(m_1,\dots,m_k)$ of several
natural numbers $m_l$. These eigenvalues are in general
degenerate with multiplicities defined essentially 
by dimensions of the representations of $G$ and $H$.
They are also polynomials\footnote{The spectrum of differential 
operators on coset spaces can be calculated exactly even if one adds
a homogeneous gauge field (see \cite{Dolan:2003bj} and references 
therein).}
 in $m_l$. The heat kernel is then
represented as an infinite sum
\begin{equation}
K(t,D)=\sum_{m_l} N(m_1,\dots,m_k) \exp (-t Q(m_1,\dots,m_k))\,.
\label{hsheatk}
\end{equation}
There exist several tricks which can be used to evaluate
the $t\to 0$ asymptotics of such sums. For example, one may use
the Poisson summation formula (\ref{eq:asPoi}). By taking derivatives
with respect to $t$ one obtains 
\begin{equation}
\sum_{l\in \mathbb{Z}} l^{2r} \exp \left( -tl^2\right)
\simeq \sqrt{\pi} 2^{-r} (2r-1)!! t^{-(2r+1)/2} +
\mathcal{O}\left( e^{-1/t} \right)\label{diffPoi}
\end{equation}
for $r\in\mathbb{N}$.

More general polynomials of $l$ may be treated by the 
Euler--Maclaurin formula which reduces sums to the integrals.
Let $F(\tau )$ be a function defined on $0 \le \tau < \infty$. If the
$2m$-th derivative
$F^{(2m)}(\tau)$ is absolutely integrable on $(0,\infty )$ 
\begin{eqnarray}
&&\sum_{i=0}^k F(i) -\int_0^\infty F(\tau )d\tau =\frac 12 [F(0)+F(k)]
\nonumber \\
&&\qquad\qquad\qquad +\sum_{s=1}^{m-1} \frac {B_{2s}}{(2s)!}
[F^{(2s-1)}(k)-F^{(2s-1)}(0)]+{\rm Rem}_m(n), \label{eq:asEM}
\end{eqnarray}
where the reminder ${\rm Rem}_m$ satisfies:
\begin{equation}
|{\rm Rem}_m(n)| \le (2-2^{1-2m}) \frac {|B_{2m}|}{(2m)!}
\int_0^n |F^{(2m)}(\tau )| d\tau .
\label{eq:asrem}
\end{equation}
$B_s$ are the Bernoulli numbers. If $F(i)$ is taken to be
$N_i \exp (-t \lambda_i )$, one can take limit $k\to \infty$
in (\ref{eq:asEM}) and restrict the summation to some finite
$m$ to obtain asymptotic series for the heat kernel.
For example, one can easily recover the expansion (\ref{diffPoi}).

In many cases calculation can be done by means of the Mellin transform.
This method will be clear from the following example.
Consider
\begin{equation}
K(t)=\sum_{k=1}^\infty k\exp (-tk^2) \label{(A.1)}
\end{equation}
The Mellin transform of $K(t)$ is
\begin{equation}
\tilde K(s)=\sum_k k \int_0^\infty dx\ x^{s-1}
\exp (-xk^2) =\sum_{k=1}^\infty k^{1-2s}\Gamma (s)
=\zeta_R(2s-1) \Gamma (s), \label{(A.2)}
\end{equation}
where $\zeta_R$ is the Riemann zeta-function. Inverse
Mellin transformation gives
\begin{equation}
K(t)=\frac 1{2\pi i} \int_C t^{-s} \zeta_R(2s-1)
\Gamma (s) ds. \label{(A.3)}
\end{equation}
The contour $C$ covers all poles of the integrand
at $s=1, 0, -1, -2,\dots$. By calculating the residues we
obtain the desired expansion:
\begin{equation}
K(t)=\frac 1{2t} -\frac 1{12} -\frac t{120}
-\frac {t^2}{504} -\frac {t^3}{1440} +O(t^4).
\end{equation}

Actually, the techniques introduced above should only be
used for complicated multi-parameter sums (see, e.g.,
\cite{Lyakhovsky:1991vv}).
For a simple one-parameter sum it is easier to
use combinations of known asymptotic series (cf. Appendix
of Ref. \cite{Birmingham:1987cy}).

In this method complexity of calculations of the heat kernel
coefficients $a_k$ is almost independent of $k$. Therefore,
the algebraic techniques were applied to higher dimensional theories
where one needs higher heat kernel coefficients to
perform renormalization or to calculate the anomalies.
Rather naturally, the most simple toroidal spaces were
considered first \cite{Appelquist:1983zs,Appelquist:1983vs,Rubin:1983zz,
Kogan:1983fp,Fradkin:1983kf,Buchbinder:1988np}, and then
computations on spheres were performed
\cite{Candelas:1984ae,Kikkawa:1985qc,Kantowski:1987ct,Birmingham:1987cy,
Birmingham:1988iv,Gleiser:1987mg,BuchOd,Buchbinder:1989cd,Elizalde:1996nb}.
Other homogeneous spaces were considered for example in 
\cite{Lyakhovsky:1991vv,Vassilevich:1991zi}. In the same approach
non-minimal operators on homogeneous spaces were treated
in \cite{Vassilevich:1992tb,Vassilevich:1993yt}. We refer to
\cite{Camporesi:1990wm} for a more extensive literature survey.
